% Part: proof-theory
% Chapter: natural-deduction
% Section: sequents

\documentclass[../../../include/open-logic-section]{subfiles}

\begin{document}

\olfileid[id]{pt}{nat}{seq}

\olsection{Deduksi Alami dengan Sekuen: \Log{N2c} dan~\Log{N2i}}

!!^a{proof}~$\delta$ dalam \Log{N1c} atau \Log{N1i} menunjukkan bahwa
!!{formula} akhirnya~$!A$ mengikuti dari asumsi-asumsi terbukanya. Jika
$\Gamma$ adalah himpunan !!{formula}~$!B$ sedemikian sehingga $!B^x$
muncul sebagai asumsi terbuka dalam~$\delta$, kita dapat menuliskannya
sebagai $\delta \Proves \Gamma \Sequent !A$. Hal ini menyarankan bahwa
kita juga dapat merumuskan sistem deduksi alami untuk sekuen. Dalam
sistem semacam itu, setiap sekuen membawa informasi tentang asumsi terbuka
mana yang menjadi tempat bergantungnya !!{formula} di sebelah kanan;
yakni, asumsi mana yang terbuka dalam sub-!!{proof} yang berakhir pada
sekuen tersebut. Inilah deduksi alami dengan sekuen, atau deduksi alami
``bergaya sekuen,'' dan sistem yang dihasilkan kita sebut
\Log{N2c} dan~\Log{N2i}.

Agar korespondensi antara \Log{N1} dan \Log{N2} lebih dekat, mari kita
ambil himpunan~$\Gamma$ di sebelah kiri sebagai himpunan yang terdiri
atas !!{formula} berlabel~$x:!B$. Karena aturan dalam \Log{N1} hanya
bekerja pada !!{formula} dalam premis, yakni !!{formula} akhir dari
sub-!!{proof} langsung, aturan \Log{N2} hanya bekerja pada sukseden
sekuen. Dengan demikian, kita dapat menyebut formula berlabel di
sebelah kiri sekuen cukup sebagai konteks.

Sebagai pengganti asumsi, bukti \Log{N2} memiliki sekuen-sekuen
teratas berupa aksioma berbentuk
\[x:!A \Sequent !A.\]
Aturannya bersesuaian tepat dengan aturan \Log{N1} dan diberikan dalam
\olref{tab:N2}.

\subfile{rules-N2}

Karena kita harus mengizinkan (seperti dalam \Log{N1c} dan \Log{N1i})
bahwa aturan \Intro{\lif}, \Intro{\lnot}, \Elim{\lor}, dan
\Elim{\lexists} boleh, tetapi tidak harus, benar-benar melepaskan
asumsi, kita juga mengizinkan penerapan yang benar atas aturan-aturan
yang bersesuaian dalam \Log{N2c} dan~\Log{N2i} ketika
!!{formula} konteks $!A$, $!B$, dan $!A(a)$ yang
bersesuaian tidak terdapat dalam konteks premis terkait.

\begin{prop}
Jika $\delta$ adalah !!a{proof} dalam~\Log{N1c} atau \Log{N1i}
atas~$!A$, maka terdapat !!a{proof}~$\delta'$ atas
$\Gamma \Sequent !A$ dalam~\Log{N2c} (atau \Log{N2i}), dengan
$\Gamma$ sebagai himpunan semua~$x:!B$ sedemikian sehingga $!B^x$
merupakan asumsi terbuka dari~$\delta$.
\end{prop}

\begin{proof}
  Melalui induksi pada $n=\pheight{\delta}$. Jika $n=0$, $\delta$
  adalah satu asumsi terbuka~$!A^x$. Maka $\Gamma=\{x:!A\}$ dan
  $\Gamma \Sequent !A$ merupakan aksioma \Log{N2c}
  atau~\Log{N2i}.

  Jika $n>0$, kita membedakan kasus menurut inferensi terakhir
  dalam~$\delta$. Kita hanya membahas satu kasus; sisanya ditinggalkan
  sebagai latihan.

  Andaikan inferensi terakhir adalah $\Intro{\lif}$. Maka $\delta$
  berakhir pada
  \[
  \AxiomC{$\Discharge{!B}{x}$}
  \RightLabel{$\delta_1$}
  \DeduceC{$!C$}
  \DischargeRule{\Intro{\lif}}{x}
  \UnaryInfC{$!B \lif !C$}
  \DisplayProof
  \]
  Berdasarkan hipotesis induksi, terdapat !!a{proof}~$\delta_1'$
  dalam~\Log{N2c} (atau \Log{N2i}) atas
  $\Gamma' \Sequent !C$, dengan $\Gamma'$ sebagai himpunan asumsi
  berlabel yang terbuka dalam~$\delta_1$. Jika $!B^x$ merupakan asumsi
  terbuka dari $\delta_1$, asumsi itu dilepaskan pada inferensi
  \Intro{\lif}, dan asumsi-asumsi terbuka dari~$\delta$ adalah
  $\Gamma=\Gamma'\setminus\{x:!B\}$. Dengan kata lain, $\delta_1'$
  merupakan !!a{proof} atas~$x:!B,\Gamma\Sequent !C$, dan kita dapat
  mengambil $\delta'$ sebagai
  \[
  \AxiomC{}
  \RightLabel{$\delta_1'$}
  \Deduce$x:!B, \Gamma \fCenter !C$
  \RightLabel{\Intro{\lif}}
  \UnaryInf$\Gamma \fCenter !B \lif !C$
  \DisplayProof
  \]
  Jika $!B^x$ bukan asumsi terbuka dari~$\delta_1$, maka
  \Intro{\lif} tidak melepaskan asumsi, dan asumsi-asumsi terbuka
  dari $\delta$ dan $\delta_1$ sama. Karena formula konteks $x:!B$
  tidak disyaratkan hadir dalam premis \Intro{\lif} pada \Log{N2c}
  dan~\Log{N2i}, kita dapat mengambil $\delta'$ sebagai
  \[
  \AxiomC{}
  \RightLabel{$\delta_1'$}
  \Deduce$\Gamma \fCenter !C$
  \RightLabel{\Intro{\lif}}
  \UnaryInf$\Gamma \fCenter !B \lif !C$
  \DisplayProof
  \]

  Kasus ketika $\delta$ berakhir pada aturan lain serupa.
\end{proof}

\begin{prop}
Jika $\delta$ adalah !!a{proof} dalam \Log{N2c} atau \Log{N2i}
atas $\Gamma \Sequent !A$, maka terdapat !!a{proof}~$\delta'$ dalam
\Log{N1c} (atau \Log{N1i}) dengan !!{formula} akhir~$!A$ dan asumsi
terbuka $!B^x$ untuk setiap $x:!B\in\Gamma$.
\end{prop}

\begin{proof}
  Kita kembali menggunakan induksi pada !!{height}~$n$ dari~$\delta$.
  Jika $n=0$, $\delta$ adalah aksioma $x:!A \Sequent !A$.
  !!{proof}~$\delta'$ adalah asumsi~$!A^x$.

  Jika $n>0$, kita membedakan kasus menurut inferensi terakhir
  dalam~$\delta$. Misalnya, jika inferensi tersebut adalah
  \Intro{\lif}, maka $\delta$ berakhir pada
  \[
  \bottomAlignProof
  \AxiomC{}
  \RightLabel{$\delta_1$}
  \Deduce$x:!B, \Gamma \fCenter !C$
  \RightLabel{\Intro{\lif}}
  \UnaryInf$\Gamma \fCenter !B \lif !C$
  \DisplayProof
  \quad\text{ atau }\quad
  \bottomAlignProof
  \AxiomC{}
  \RightLabel{$\delta_1$}
  \Deduce$\Gamma \fCenter !C$
  \RightLabel{\Intro{\lif}}
  \UnaryInf$\Gamma \fCenter !B \lif !C$
  \DisplayProof
  \]
  Berdasarkan hipotesis induksi, terdapat !!a{proof}~$\delta_1'$
  dalam~\Log{N1c} (atau \Log{N1i}) atas~$!C$. Dalam kasus pertama,
  $!B^x$ merupakan asumsi terbuka dari~$\delta_1'$; dalam kasus kedua
  tidak. Kita dapat mengambil !!{proof}~$\delta'$ sebagai
  \[
  \bottomAlignProof
  \AxiomC{$\Discharge{!B}{x}$}
  \RightLabel{$\delta_1'$}
  \DeduceC{$!C$}
  \DischargeRule{\Intro{\lif}}{x}
  \UnaryInfC{$!B \lif !C$}
  \DisplayProof
  \quad\text{ atau }\quad
  \bottomAlignProof
  \AxiomC{}
  \RightLabel{$\delta_1'$}
  \DeduceC{$!C$}
  \RightLabel{\Intro{\lif}}
  \UnaryInfC{$!B \lif !C$}
  \DisplayProof
  \]
  secara berurutan.
\end{proof}

\end{document}
